% theory/bias_law/notes/polar_bias.tex -- research note extending "From SAGA to Muon-type steps" (paper/paper_saga_to_muon.tex).
% Build (from theory/bias_law/notes/):  pdflatex polar_bias && pdflatex polar_bias
% Numbers: experiments/bias_law/checks/c01..c07 (logs in results/bias_law/checks/).  Literature: theory/bias_law/notes/novelty_ledger.md, bibliography/bias_law_literature_1.md.
\documentclass[10pt]{article}
\usepackage[utf8]{inputenc}
\usepackage[T1]{fontenc}
\usepackage{amsmath,amssymb,amsthm,mathtools}
\usepackage{booktabs}
\usepackage{enumitem}
\usepackage[margin=1in]{geometry}
\usepackage[colorlinks=true,linkcolor=blue!55!black,citecolor=blue!55!black,urlcolor=blue!55!black]{hyperref}
\usepackage{xcolor}
\setlength{\emergencystretch}{2em}

\newtheorem{theorem}{Theorem}
\newtheorem{proposition}[theorem]{Proposition}
\newtheorem{lemma}[theorem]{Lemma}
\newtheorem{corollary}[theorem]{Corollary}
\theoremstyle{definition}
\newtheorem{assumption}{Assumption}
\renewcommand{\theassumption}{\Alph{assumption}}
\theoremstyle{remark}
\newtheorem{remark}[theorem]{Remark}

\newcommand{\R}{\mathbb{R}}
\newcommand{\E}{\mathbb{E}}
\newcommand{\Prob}{\mathbb{P}}
\newcommand{\nop}[1]{\|#1\|_{\mathrm{op}}}
\newcommand{\nnuc}[1]{\|#1\|_{*}}
\newcommand{\nF}[1]{\|#1\|_{F}}
\newcommand{\inner}[2]{\langle #1, #2\rangle}
\newcommand{\tr}{\operatorname{tr}}
\newcommand{\polar}{\operatorname{polar}}
\newcommand{\Med}{\operatorname{Med}}
\newcommand{\He}{\mathrm{He}}
\newcommand{\cA}{\mathsf{A}}
\newcommand{\cB}{\mathsf{B}}
\newcommand{\ones}{\mathbf{1}}
\newcommand{\Lop}{L_{\mathrm{op}}}
\newcommand{\muop}{\mu_{\mathrm{op}}}
\newcommand{\cF}{\mathcal{F}}
\newcommand{\st}{\dagger}

\title{\vspace{-2em}The first-order bias law of stochastic orthogonalized steps:\\
two cubic moments, an exact dimension factor, and momentum as a batch multiplier}
\author{Research note extending ``From SAGA to Muon-type steps'' (folder \texttt{*/bias\_law/})}
\date{25 November 2026}

\begin{document}
\maketitle

\begin{abstract}
A stochastic polar step $W^+=W-\eta\,\polar(\nabla F(W)+\xi)$ does not converge to the minimizer when the gradient
noise is skewed: the manuscript shows that the limit gradient is minus a nuclear median of the noise and gives a
\emph{formal} $1/b$ expansion for i.i.d.\ entries, with a constant $a_{K,d}/c_{K,d}$ known only by Monte Carlo. We
turn this into a theory. (1)~\emph{A Cornish--Fisher law for norm medians}, rigorous under Cram\'er's condition, for
every noise law and every norm whose subgradient map satisfies an oscillation condition (verified for the sign,
Euclidean and spectral geometries): the minimizers of the noise-smoothed dual norm lie within $o(1/b)$ of
$-\frac1{6b}J_\Sigma^{-1}T_\Sigma[K_3]$, where $J_\Sigma,T_\Sigma$ are the second and fourth derivatives at the origin
of the Gaussian-smoothed dual norm and $K_3$ is the third-cumulant tensor. (2)~\emph{The two-moment theorem.} For
isotropic covariance and the spectral geometry, orthogonal equivariance collapses the dependence on the
$\binom{Kd+2}3$-dimensional tensor $K_3$ to two cubic matrix moments:
$G^\st_b=\frac1{6\sigma^2b}\big(a_1\E[\zeta\nF\zeta^2]+a_2\E[\zeta\zeta^\top\zeta]\big)+o(1/b)$, with $a_1,a_2$
rational functions of $K,d$ and of one number $q_{K,d}=\E\sum_{i\ne j}\frac{s_is_j}{s_i+s_j}/\E\sum_is_i$ of a Gaussian
matrix. Consequences are exact: the constant of the manuscript is
$a_{K,d}/c_{K,d}=\frac{3(3+2q_{K,d})}{(K+2)(d+2)}\in\big[\frac{9}{(K+2)(d+2)},\frac3{\max(K,d)+2}\big]$, it is
$\sim\theta(\gamma)/d$ with $\theta(\gamma)\in[\frac{12}5,3]$ in the proportional limit, $a_1<0<a_2$ on every size
computed, the fourth derivative of the smoothed nuclear norm is negative in every direction, and neither sign descent
nor orthogonalization dominates the other. (3)~\emph{State-dependent noise}: all zeros of the mean field near a
nondegenerate critical point are at $W_\ast+\frac1bH^{-1}g_\ast+o(1/b)$; evaluated on the finite-sum
problem of the manuscript (discrete noise, where the law is not proved), the formula gives $0.314$ with no fitted
parameter for a quantity the manuscript measures rising to $0.308$ at $b=32$. (4)~\emph{Momentum}: polar steps on an
exponential moving average with Robbins--Monro lengths converge almost surely to minus a nuclear median of the
stationary filtered noise, whose first-order bias is that of batch size $b\,\frac{1+\beta+\beta^2}{1-\beta^2}$, not the
variance-based $b\,\frac{1+\beta}{1-\beta}$. (5)~Under sampling without replacement the nuclear medians obey an exact
complement duality, $0$ is a median at $b=n/2$, and the offset changes sign beyond (first-order factor formal, sign
change observed numerically).
\end{abstract}

\section{Introduction}

\paragraph{Starting point.} For a random matrix $\xi\in\R^{K\times d}$ the manuscript \cite{manuscript} defines the
noise-smoothed nuclear norm $\psi_\xi(G)=\E[\nnuc{G+\xi}-\nnuc\xi]$, the mean field $\Phi_\xi(G)=\E\polar(G+\xi)\in
\partial\psi_\xi(G)$ and the set of nuclear medians $\Med_\ast(\xi)=-\arg\min\psi_\xi$, and proves
(\cite[Theorem~2]{manuscript}) that for $0<\muop\le\Lop<\infty$, i.i.d.\ additive noise and deterministic lengths
with $\sum_t\eta_t=\infty$, $\sum_t\eta_t^2<\infty$, the gradients $\nabla F(W_t)$ of
$W_{t+1}=W_t-\eta_t\polar(\nabla F(W_t)+\xi_t)$ converge almost surely to a point of $-\Med_\ast(\xi)$. For the mean
$\xi_b$ of $b$ i.i.d.\ matrices with i.i.d.\ entries, a formal Edgeworth computation
(\cite[Proposition~4]{manuscript}) gives the minimizer $G_b^\st=\frac{\kappa_3}{6\sigma^2b}\frac{a_{K,d}}{c_{K,d}}\ones\ones^\top+o(1/b)$;
the archived notes of the project also contain the formal operator form $G^\st_b=-\frac1b\mathcal J^{-1}\mathcal T$ for
a general covariance and, in the scalar case with a frozen iterate, the momentum factor
$\frac{1-\beta^2}{1+\beta+\beta^2}$. None of this is proved there; the constant $a_{K,d}/c_{K,d}$ is tabulated by
Monte Carlo (``$d\,a/c$ between $1$ and $2.8$''); momentum is excluded from the convergence theorem; and the noise is
additive. These are our starting point, not our results.

\paragraph{Results.} Section~\ref{sec:cf} proves the expansion for every noise law satisfying Cram\'er's condition and
every norm satisfying Assumption~\ref{ass:O} (Theorem~\ref{thm:cf}). Section~\ref{sec:two} is the core of the note: the two-moment theorem
(Theorem~\ref{thm:two}), which makes the bias operator explicit for the spectral geometry and yields the exact dimension
dependence, sharp bounds, the proportional limit, a sign structure and a comparison with sign descent and normalized
steps (Corollaries~\ref{cor:dim}--\ref{cor:compare}). Section~\ref{sec:state} treats parameter-dependent noise
(Theorem~\ref{thm:state}), Section~\ref{sec:mom} momentum (Theorems~\ref{thm:momconv} and~\ref{thm:mombias}),
Section~\ref{sec:wor} sampling without replacement (Proposition~\ref{prop:wor}). Section~\ref{sec:num} reports the
numerical checks and Section~\ref{sec:ledger} the novelty ledger and the limits of the note.

\paragraph{Status convention.} Every theorem below is proved in full, up to the quoted results that are named at the
place of use: Theorem~20.1 of Bhattacharya and Rao \cite{br}, the perturbation bound of Li and Sun \cite{lisun} for the
polar factor, the divergence formula for spectral functions \cite{candes}, and \cite[Theorem~2]{manuscript}. We did not
have the book \cite{br} at hand: the statement used is the i.i.d.\ case of the one quoted verbatim in
\cite[Theorem~4.1]{angstpoly} and should be checked against the book before submission. Statements that are only
formal or numerical are labelled as such.

\section{Setting}\label{sec:setting}

$E=\R^{K\times d}$ with the Frobenius inner product, $N=Kd$, $r=\min(K,d)$. A \emph{geometry} is a norm $\nu$ on $E$
(the dual of the norm in which steps are measured) with a Borel map $o:E\to E$ such that $o(x)\in\partial\nu(x)$,
$o(tx)=o(x)$ for $t>0$ and $o(-x)=-o(x)$; then $|o|\le C_o$ for a constant. The three examples are the \emph{spectral}
geometry ($\nu=\nnuc\cdot$, $o=\polar$: Muon-type steps), the \emph{sign} geometry ($\nu=\|\cdot\|_1$, $o=\operatorname{sign}$)
and the \emph{Euclidean} one ($\nu=\nF\cdot$, $o(x)=x/\nF x$: normalized steps). For a random $\xi\in E$ with
$\E\nu(\xi)<\infty$,
\[
\psi_\xi(G)=\E[\nu(G+\xi)-\nu(\xi)],\qquad \Phi_\xi(G)=\E\,o(G+\xi)\in\partial\psi_\xi(G),\qquad
\Med_\nu(\xi)=-\arg\min\psi_\xi .
\]
$\psi_\xi$ is convex, finite, and $\psi_\xi(G)\ge\nu(G)-2\E\nu(\xi)$, so $\arg\min\psi_\xi$ is nonempty, convex and compact.

\paragraph{Gaussian objects.} For a positive definite covariance $\Sigma$ on $E$ and $Z_\Sigma\sim N(0,\Sigma)$ let
$\psi_\Sigma=\psi_{Z_\Sigma}$, $\Phi_\Sigma=\Phi_{Z_\Sigma}$, and
\begin{equation}\label{eq:JT}
J_\Sigma=D\Phi_\Sigma(0)=\nabla^2\psi_\Sigma(0)=\E\big[o(Z_\Sigma)\otimes\Sigma^{-1}Z_\Sigma\big],\qquad
T_\Sigma[K_3]=D^3\Phi_\Sigma(0)[K_3]=\E\big[o(Z_\Sigma)\,\inner{K_3}{\He_3^\Sigma(Z_\Sigma)}\big],
\end{equation}
for a symmetric $3$-tensor $K_3$ on $E$, where $\He_3^\Sigma(x)=y^{\otimes3}-3\,\mathrm{sym}(y\otimes\Sigma^{-1})$,
$y=\Sigma^{-1}x$, and $D^3\Phi(h)[K_3]=\sum_{abc}K_3^{abc}\partial_a\partial_b\partial_c\Phi(h)$. In other words
$T_\Sigma[K_3]=\nabla^4\psi_\Sigma(0)[K_3,\cdot\,]$.

\begin{lemma}[Gaussian side]\label{lem:gauss}
$\Phi_\Sigma=\nabla\psi_\Sigma$ is odd, $C^\infty$, with derivatives of every order bounded on $E$, uniformly for
$\Sigma$ in a compact set of positive definite matrices; the identities \eqref{eq:JT} hold; $D^2\Phi_\Sigma(0)=0$;
$J_\Sigma$ is symmetric positive definite; $J_{c\Sigma}=c^{-1/2}J_\Sigma$ and $T_{c\Sigma}=c^{-3/2}T_\Sigma$ for $c>0$;
and there is $c_0>0$ with $|\Phi_\Sigma(h)|\ge c_0\min(|h|,1)$.
\end{lemma}
\begin{proof}
$\Phi_\Sigma(h)=\int o(y)\varphi_\Sigma(y-h)\,dy$ with $o$ bounded: differentiate under the integral; the $k$-th
derivative is bounded by $C_o\int|\partial^k\varphi_\Sigma|$, and $(-1)^k\partial^k\varphi_\Sigma=\He_k^\Sigma\varphi_\Sigma$
gives \eqref{eq:JT}. $\nu$ is differentiable almost everywhere with gradient $o$, so $\nabla\psi_\Sigma=\Phi_\Sigma$ by
dominated convergence. Oddness of $o$ and symmetry of $Z_\Sigma$ give the oddness of $\Phi_\Sigma$, hence
$D^2\Phi_\Sigma(0)=0$. The scalings follow from $o(tx)=o(x)$. Positivity: fix $v\ne0$, let
$s=\inner{\Sigma^{-1}v}{Z}/\inner v{\Sigma^{-1}v}$ and $Z'=Z-sv$, which are independent, $s\sim N(0,\tau^2)$. Then
$\psi_\Sigma(tv)+\E\nu(Z)=\E\,h(s+t)$ with $h(\rho)=\E\nu(Z'+\rho v)$ convex and
$|\rho|\nu(v)-\E\nu(Z')\le h(\rho)\le\E\nu(Z')+|\rho|\nu(v)$, so $h'$ increases from $-\nu(v)$ to $\nu(v)$ and
$\inner{J_\Sigma v}v=\int\varphi_\tau\,dh'>0$. Finally $t\mapsto\inner{\Phi_\Sigma(t\hat h)}{\hat h}$ is nondecreasing
(convexity), equals $t\inner{J_\Sigma\hat h}{\hat h}+O(t^3)$ near $0$, and the constants are uniform on the unit sphere.
\end{proof}

\begin{assumption}[noise]\label{ass:N}
$\xi_b=\frac1b\sum_{i\le b}\zeta_i$ with $\zeta_i$ i.i.d.\ in $E$, $\E\zeta=0$, $\E|\zeta|^3<\infty$, covariance
$\Sigma\succ0$, third cumulant tensor $K_3=\E\,\zeta^{\otimes3}$, and Cram\'er's condition
$\limsup_{|t|\to\infty}|\E e^{i\inner t\zeta}|<1$.
\end{assumption}
\begin{assumption}[oscillation]\label{ass:O}
For every compact set $\mathcal K$ of positive definite covariances there are $C,a>0$ with
$\sup_{\Sigma\in\mathcal K}\sup_{y\in E}\E\,\mathrm{osc}\big(o;B(y+Z_\Sigma,\varepsilon)\big)\le C\varepsilon^a$ for
$\varepsilon\le1$, where $\mathrm{osc}(o;B)=\max_{k}(\sup_Bo_k-\inf_Bo_k)$ over the coordinates $o_k$.
\end{assumption}

\begin{lemma}\label{lem:osc}
Assumption~\ref{ass:O} holds for the sign, Euclidean and spectral geometries (with any $a<1$). Moreover, for
$Z_\Sigma$ as above, $\sup_y\Prob(\sigma_r(y+Z_\Sigma)\le t)\le C_\Sigma t$.
\end{lemma}
\begin{proof}
Sign: the oscillation vanishes unless some coordinate of $y+Z$ is within $\varepsilon$ of $0$, an event of probability at
most $N\varepsilon\sup\varphi$. Euclidean: $\mathrm{osc}\le\min(2,2\varepsilon/(|x|-\varepsilon)_+)$. Spectral, $K\le d$
(transpose otherwise): with rows $\rho_j$ of $X$, $\sigma_K(X)\ge K^{-1/2}\min_j\operatorname{dist}(\rho_j,\operatorname{span}\{\rho_l\}_{l\ne j})$;
conditionally on the other rows, $\rho_j$ is Gaussian with covariance at least $\lambda_{\min}(\Sigma)I$ and its
component along a unit normal to the span has a density bounded by $(2\pi\lambda_{\min})^{-1/2}$; a union bound gives the
small-ball estimate. By \cite{lisun}, $\nF{\polar(A)-\polar(B)}\le2\nF{A-B}/(\sigma_r(A)+\sigma_r(B))$ for $A,B$ of rank
$r$; hence $\mathrm{osc}(\polar;B(x,\varepsilon))\le\min(2,4\varepsilon/\sigma_r(x))$ when $\sigma_r(x)>2\varepsilon$, and
$\E\min(2,4\varepsilon/\sigma_r)\le2\int_0^1\min(1,C\varepsilon/u)\,du\le C'\varepsilon\log(e/\varepsilon)$.
\end{proof}

\section{A Cornish--Fisher law for norm medians}\label{sec:cf}

\begin{lemma}[Edgeworth expansion of the mean field, uniform in the shift]\label{lem:edge}
Under Assumptions~\ref{ass:N} and~\ref{ass:O}, with $X_b=\sqrt b\,\xi_b$,
\[
\sup_{h\in E}\Big|\E\,o(X_b+h)-\Phi_\Sigma(h)-\frac1{6\sqrt b}\,D^3\Phi_\Sigma(h)[K_3]\Big|=o(b^{-1/2}).
\]
\end{lemma}
\begin{proof}
Theorem~20.1 of \cite{br} with $s=3$: there are $\delta_b=o(b^{-1/2})$ and $c,c'>0$, depending on the law of $\zeta$
only, such that for every bounded Borel $f$,
$|\int f\,d(Q_b-\Phi_{0,\Sigma}-b^{-1/2}P_1)|\le\sup|f|\,\delta_b+c'\int\mathrm{osc}(f;B(x,2e^{-cb}))\varphi_\Sigma(x)dx$,
where $Q_b$ is the law of $X_b$ and $P_1$ has density $-\frac16\sum K_3^{abc}\partial_{abc}\varphi_\Sigma$. Apply it to
the coordinates of $f_h=o(\cdot+h)$: the supremum is at most $C_o$, the oscillation term is at most
$C(2e^{-cb})^a$ by Assumption~\ref{ass:O}, uniformly in $h$, and
$\int o(x+h)(-\tfrac16)K_3^{abc}\partial_{abc}\varphi_\Sigma(x)dx=\frac16K_3^{abc}\partial_{h_a}\partial_{h_b}\partial_{h_c}\int
o(y)\varphi_\Sigma(y-h)dy$.
\end{proof}

\begin{theorem}[first-order law, any norm, any noise]\label{thm:cf}
Under Assumptions~\ref{ass:N} and~\ref{ass:O} let
\begin{equation}\label{eq:gstar}
g_\ast=-\tfrac16\,J_\Sigma^{-1}\,T_\Sigma[K_3]\ \in E .
\end{equation}
For every $\varepsilon>0$ there is $b_0$ such that for all $b\ge b_0$
\[
\arg\min\psi_{\xi_b}\ \subset\ \Big\{G:\ \Big|G-\frac{g_\ast}b\Big|\le\frac\varepsilon b\Big\},\qquad\text{i.e.}\qquad
\Med_\nu(\xi_b)=-\frac{g_\ast}b+o\Big(\frac1b\Big).
\]
In the spectral geometry, the limit gradient of the polar iteration of \cite[Theorem~2]{manuscript} with noise
$\xi_b$ is therefore $g_\ast/b+o(1/b)$, and if $F$ is $C^2$ at $W_\ast$ with Hessian $H$ the limit point $W^\st_b$ has
$W^\st_b-W_\ast=\frac1bH^{-1}g_\ast+o(\frac1b)$ and $F(W^\st_b)-F_\ast=\frac1{2b^2}\inner{g_\ast}{H^{-1}g_\ast}+o(b^{-2})$.
\end{theorem}
\begin{proof}
By homogeneity $\Phi_{\xi_b}(G)=\E\,o(X_b+\sqrt b\,G)$. Let $G=(g_\ast+u)/b$ with $|u|\le R$, so $h=\sqrt bG=O(b^{-1/2})$.
By Lemma~\ref{lem:gauss}, $\Phi_\Sigma(h)=J_\Sigma h+O(|h|^3)$ and $D^3\Phi_\Sigma(h)[K_3]=T_\Sigma[K_3]+O(|h|)$; with
Lemma~\ref{lem:edge},
\[
\Phi_{\xi_b}(G)=b^{-1/2}\big(J_\Sigma(g_\ast+u)+\tfrac16T_\Sigma[K_3]\big)+o(b^{-1/2})=b^{-1/2}J_\Sigma u+o(b^{-1/2}),
\]
uniformly in $|u|\le R$. Take $|u|=\varepsilon$: $\inner{\Phi_{\xi_b}(G)}{u}\ge b^{-1/2}(\lambda_{\min}(J_\Sigma)\varepsilon^2-o(1))>0$
for $b\ge b_0$. If a minimizer $G'$ of $\psi_{\xi_b}$ had $|G'-g_\ast/b|>\varepsilon/b$, let $G$ be the point of the
segment $[g_\ast/b,G']$ on the sphere of radius $\varepsilon/b$, so $G'-G=\lambda u$ with $\lambda>0$; since
$\Phi_{\xi_b}(G)\in\partial\psi_{\xi_b}(G)$, $\psi_{\xi_b}(G')\ge\psi_{\xi_b}(G)+\lambda\inner{\Phi_{\xi_b}(G)}u>\psi_{\xi_b}(G)$,
a contradiction. The last sentence is \cite[Theorem~2]{manuscript} and a Taylor expansion.
\end{proof}

\begin{remark}[what the theorem contains]\label{rem:cf}
(i) Sign geometry, diagonal $\Sigma$: $J=\operatorname{diag}(\sqrt{2/\pi}/\sigma_k)$,
$T[K_3]_k=-\sqrt{2/\pi}\,\kappa_{3,k}/\sigma_k^3$, so $g_{\ast,k}=\kappa_{3,k}/(6\sigma_k^2)$: the Cornish--Fisher median of a
mean \cite{cornish}, coordinate by coordinate. (ii) Euclidean geometry, $\Sigma=\sigma^2I$ on $\R^N$:
$g_\ast=\E[\zeta|\zeta|^2]/(2\sigma^2(N+2))$ (Theorem~\ref{thm:two} with $K=1$). The geometric median of the law of a
sample mean is the object of \cite{minsker}, who bound its distance to the mean by $O(b^{-1/2})$; Schwank and Drton
\cite{schwank} prove $m_i=\mu_i-\alpha_i/p+o(1/p)$ as the \emph{dimension} $p\to\infty$ at fixed $b$, with
$\alpha_i=\frac1{2\bar\sigma^2}\sum_j\E[X_iX_j^2]$: the same driving moment, and our fixed-dimension constant has $p+2$ in
place of $p$. (iii) Cram\'er's condition excludes lattice laws, for which the median of a mean oscillates
(\cite{manuscript}, archived notes). For discrete non-lattice laws (minibatches of a finite data set, drawn with
replacement) one-term expansions with remainder $O(b^{-1})$ are available under weak Cram\'er conditions and four
moments \cite{angstpoly,song}; we have not checked their hypotheses on a finite sum and state the finite-sum case as
numerical (Section~\ref{sec:num}).
\end{remark}

\section{The two-moment theorem}\label{sec:two}

From now on the geometry is spectral, $o=\polar$, and $Z$ is a standard Gaussian $K\times d$ matrix with singular values
$s_1,\dots,s_r$. For a symmetric $3$-tensor $K_3$ on $E$ define the two contractions
\[
\cA[K_3]_{ij}=\sum_{k,l}K_3^{(ij),(kl),(kl)},\qquad \cB[K_3]_{ij}=\sum_{k,l}K_3^{(kj),(kl),(il)};
\]
for $K_3=\E\,\zeta^{\otimes3}$ these are $\cA=\E[\zeta\nF\zeta^2]$ and $\cB=\E[\zeta\zeta^\top\zeta]$. Let
\begin{equation}\label{eq:q}
m_p=\E\sum_is_i^p,\qquad q_{K,d}=\frac1{m_1}\,\E\sum_{i\ne j}\frac{s_is_j}{s_i+s_j},\qquad c_{K,d}=\frac{m_1}{Kd}.
\end{equation}

\begin{theorem}[two-moment theorem]\label{thm:two}
Let $K,d\ge2$, let Assumption~\ref{ass:N} hold with $\Sigma=\sigma^2I$, and let $D_{K,d}=(K-1)(d-1)(K+2)(d+2)$,
\begin{equation}\label{eq:a12}
a_1=\frac{3\big[(K-1)(d-1)-2(K+d+1)\,q_{K,d}\big]}{D_{K,d}},\qquad
a_2=\frac{6\big[(Kd+2)\,q_{K,d}+(K-1)(d-1)\big]}{D_{K,d}} .
\end{equation}
\begin{enumerate}[label=(\alph*),itemsep=1pt,topsep=2pt]
\item \emph{(Reduction.)} $T_I[K_3]=-c_{K,d}\,(a_1\cA[K_3]+a_2\cB[K_3])$ and $J_I=c_{K,d}\,\mathrm{Id}$: the nuclear medians
of $\xi_b$ satisfy
\begin{equation}\label{eq:two}
-\Med_\ast(\xi_b)=\frac1{6\sigma^2b}\Big(a_1\,\E\big[\zeta\nF\zeta^2\big]+a_2\,\E\big[\zeta\zeta^\top\zeta\big]\Big)+o\Big(\frac1b\Big).
\end{equation}
The bias operator $K_3\mapsto g_\ast$ has rank $Kd$ on a space of dimension $\binom{Kd+2}3$.
\item \emph{(One number.)} $m_3=(K+d-q_{K,d})\,m_1$; $0<q_{K,d}<\frac{r-1}2$; $a_2>0$.
\item \emph{(Rank-one and full-rank directions.)}
$a_1+a_2=\dfrac{3(3+2q_{K,d})}{(K+2)(d+2)}$ and, with $\bar d=\max(K,d)$,
$r\,a_1+a_2=\dfrac{3(\bar d-1-2q_{K,d})}{(\bar d-1)(\bar d+2)}>0$.
\item \emph{(Negative fourth derivative.)} For every $D\ne0$,
$\nabla^4\psi_I(0)[D,D,D,D]=-c_{K,d}\big(a_1\nF D^4+a_2\tr(DD^\top)^2\big)<0$.
\item \emph{(Proportional limit.)} If $K,d\to\infty$ with $K/d\to\gamma\in(0,1]$, then $q_{K,d}/K\to\kappa(\gamma)$,
$d\,a_2\to6\kappa(\gamma)$, $d^2a_1\to\frac3\gamma\big(1-2(1+\gamma)\kappa(\gamma)\big)$, with
\[
\kappa(\gamma)=\frac{(1+\gamma)\mu_{1/2}(\gamma)-\mu_{3/2}(\gamma)}{\gamma\,\mu_{1/2}(\gamma)},\qquad
\mu_s(\gamma)=\int x^s\,d\mathrm{MP}_\gamma(x)={}_2F_1(-s,1-s;2;\gamma),
\]
$\kappa(1)=\frac25$ and $\kappa(0^+)=\frac12$.
\end{enumerate}
For $K=1$ (or $d=1$) the two contractions coincide and
$-\Med_\ast(\xi_b)=\frac{3}{(d+2)}\frac{\E[\zeta|\zeta|^2]}{6\sigma^2b}+o(1/b)$.
\end{theorem}

\begin{proof}
\emph{(a) Reduction.} $J_I$ and $T_I$ are equivariant under $Z\mapsto UZV^\top$, $U\in O(K)$, $V\in O(d)$, because
$\polar(UZV^\top)=U\polar(Z)V^\top$ and the law of $Z$ is invariant. $E=\R^K\otimes\R^d$ is irreducible, so
$J_I=c\,\mathrm{Id}$ with $c=\tr J_I/N=\E\inner{\polar(Z)}Z/N=c_{K,d}$ \cite{xie,chen}. $T_I$ is an invariant element of
$E^{\otimes4}=(\R^K)^{\otimes4}\otimes(\R^d)^{\otimes4}$; the invariants of $O(K)\times O(d)$ are spanned by the
products $\pi\otimes\pi'$ of an $O(K)$-invariant and an $O(d)$-invariant, and by the first fundamental theorem for the
orthogonal group each factor is spanned by the three pairings of the four indices. Label a pairing by the input slot
paired with the output index; $T_I$ is symmetric in its three input slots, the symmetric group $S_3$ permutes the nine
products $(\pi,\pi')$ with two orbits, $\{\pi=\pi'\}$ and $\{\pi\ne\pi'\}$, and the orbit averages are the tensors of
$\cA$ and $\cB$ on symmetric arguments. Hence $T_I=\alpha\cA+\beta\cB$. To find $\alpha,\beta$ contract with one representative of each orbit:
using $\He_3(Z)^{(ab),(cd),(ef)}=Z_{ab}Z_{cd}Z_{ef}-Z_{ab}\delta_{(cd),(ef)}-Z_{cd}\delta_{(ab),(ef)}-Z_{ef}\delta_{(ab),(cd)}$,
\begin{align*}
\textstyle\sum_{ij}\cA\big[\E\,\polar(Z)_{ij}\He_3(Z)\big]_{ij}&=\E\big[\nnuc Z(\nF Z^2-N-2)\big]=-m_1,\\
\textstyle\sum_{ij}\cB\big[\E\,\polar(Z)_{ij}\He_3(Z)\big]_{ij}&=\E\big[\inner{\polar(Z)}{ZZ^\top Z}\big]-(K+d+1)m_1=m_3-(K+d+1)m_1,
\end{align*}
where $\E[\nnuc Z\nF Z^2]=(N+1)m_1$ because $\nF Z$ is independent of $Z/\nF Z$ and $\E\chi_N^3=(N+1)\E\chi_N$. On the
other side, by Lemma~\ref{lem:count}, the same contractions of $\alpha\cA+\beta\cB$ are
$\frac N3[(N+2)\alpha+(K+d+1)\beta]$ and $\frac N3[(K+d+1)\alpha+\frac12(N+K+d+3)\beta]$. The determinant of this system
is $\frac{N^2}{18}(K-1)(d-1)(K+2)(d+2)$; solving, and substituting $m_3=(K+d-q)m_1$ from (b), gives
$\alpha=-c_{K,d}a_1$, $\beta=-c_{K,d}a_2$ with \eqref{eq:a12}. Then \eqref{eq:two} is Theorem~\ref{thm:cf} with the scaling
of Lemma~\ref{lem:gauss}. The map $K_3\mapsto(\cA[K_3],\cB[K_3])$ is onto $E\oplus E$ (the two equivariant maps are
linearly independent and $E$ is irreducible), so $K_3\mapsto a_1\cA+a_2\cB$ is onto $E$.

\emph{(b)} $\sum_is_i^3=\inner Z{f(Z)}$ with the spectral function $f(Z)=(ZZ^\top)^{1/2}Z=U\operatorname{diag}(s_i^2)V^\top$,
which is $C^1$; by Stein's identity $m_3=\E\operatorname{div}f(Z)$ and by the divergence formula for spectral functions
\cite{candes}, with $\varphi(s)=s^2$,
$\operatorname{div}f=\sum_i\big(\varphi'(s_i)+|K-d|\frac{\varphi(s_i)}{s_i}\big)+\sum_{i\ne j}\frac{s_i\varphi(s_i)-s_j\varphi(s_j)}{s_i^2-s_j^2}
=(K+d)\sum_is_i-\sum_{i\ne j}\frac{s_is_j}{s_i+s_j}$, since
$\frac{s_i^3-s_j^3}{s_i^2-s_j^2}=s_i+s_j-\frac{s_is_j}{s_i+s_j}$. The bounds on $q$ follow from
$0<\frac{s_is_j}{s_i+s_j}<\frac{s_i+s_j}4$ almost surely, and $a_2>0$ from \eqref{eq:a12}.

\emph{(c)} Algebra from \eqref{eq:a12}, using $Kd+1-(K+d)=(K-1)(d-1)$ and, for $K\le d$,
$Kd+2-K(K+d+1)=-(K-1)(K+2)$; positivity from $q<\frac{r-1}2\le\frac{\bar d-1}2$.

\emph{(d)} $\nabla^4\psi_I(0)[D,D,D,D]=\inner D{T_I[D^{\otimes3}]}$, $\cA[D^{\otimes3}]=\nF D^2D$,
$\cB[D^{\otimes3}]=DD^\top D$. If $a_1\ge0$ the sign is clear. If $a_1<0$, $\nF D^4\le\operatorname{rank}(D)\tr(DD^\top)^2$
gives $a_1\nF D^4+a_2\tr(DD^\top)^2\ge(r\,a_1+a_2)\tr(DD^\top)^2>0$ by (c).

\emph{(e)} The empirical distribution of the $s_i^2/d$ converges almost surely to the Marchenko--Pastur law of ratio
$\gamma$, with convergence of the moments $m_1/(K\sqrt d)$ and $m_3/(Kd^{3/2})$ (the extreme singular values are bounded in
every $L^p$); by (b), $q/K=(K+d)/K-m_3/(Km_1)\to[(1+\gamma)\mu_{1/2}-\mu_{3/2}]/(\gamma\mu_{1/2})$. The hypergeometric
form of $\mu_s$ is the Narayana expansion for integer $s$; for $s=\frac12,\frac32$ we checked it numerically to eight
digits and do not prove it (only the integral definition of $\mu_s$ is used); $\mu_{1/2}(1)=\frac8{3\pi}$, $\mu_{3/2}(1)=\frac{64}{15\pi}$
(quarter-circle law), hence $\kappa(1)=2-\frac85$. The limits of $a_1,a_2$ follow from \eqref{eq:a12}.

\emph{$K=1$.} The same computation with a single invariant; $m_3=(d+1)m_1$ and the contraction of the invariant with
itself is $N(N+2)/3$, so $\alpha+\beta=-3m_1/(d(d+2))$.
\end{proof}

\begin{lemma}[counting]\label{lem:count}
Let $\mathsf A_1$, $\mathsf B_1$ be one pairing tensor from each orbit and $\bar{\mathsf A}$, $\bar{\mathsf B}$ the orbit
averages (so that $\bar{\mathsf A}$, $\bar{\mathsf B}$ are the tensors of $\cA$, $\cB$ on symmetric arguments). Then
$\inner{\mathsf A_1}{\bar{\mathsf A}}=\frac{N(N+2)}3$, $\inner{\mathsf A_1}{\bar{\mathsf B}}=\inner{\mathsf B_1}{\bar{\mathsf A}}=\frac{N(K+d+1)}3$,
$\inner{\mathsf B_1}{\bar{\mathsf B}}=\frac{N(N+K+d+3)}6$.
\end{lemma}
\begin{proof}
The inner product of two pairing tensors is $K^{\ell}d^{\ell'}$, $\ell,\ell'\in\{1,2\}$ being the numbers of closed
loops of row and of column indices. The three inner products of $\mathsf A_1$ with the first orbit are $N^2,N,N$; its
six inner products with the second orbit are $Kd^2,Kd^2,K^2d,K^2d,Kd,Kd$; the six inner products of $\mathsf B_1$ with
its orbit are $K^2d^2,K^2d,Kd^2,Kd,Kd,Kd$. (Enumeration; checked by brute force for seven sizes in
\texttt{checks/c01\_two\_constants.py}.)
\end{proof}

\begin{corollary}[exact dimension factor; a question of the manuscript]\label{cor:dim}
For i.i.d.\ entries of variance $\sigma^2$ and third cumulant $\kappa_3$ (the case of \cite[Proposition~4]{manuscript}),
$\cA=\cB=\kappa_3\ones\ones^\top$ and the constant $a_{K,d}/c_{K,d}$ of the manuscript is
\[
\frac{a_{K,d}}{c_{K,d}}=a_1+a_2=\frac{3\,(3+2q_{K,d})}{(K+2)(d+2)},\qquad
\frac{9}{(K+2)(d+2)}<\frac{a_{K,d}}{c_{K,d}}<\frac{3}{\max(K,d)+2}\quad(K,d\ge2),
\]
with equality to $3/(d+2)$ for $K=1$, and $\max(K,d)\,a_{K,d}/c_{K,d}\to\theta(\gamma)=6\kappa(\gamma)$ in
the proportional limit, with $\theta(1)=\frac{12}5$ for square matrices, $\theta(0^+)=3$, $\theta\le3$ by the bound above,
and $\theta$ decreasing on $(0,1]$ numerically (not proved). Per entry, the first-order offset of a polar step is
therefore at most $3/(\max(K,d)+2)$ times that of sign descent, for every size.
\end{corollary}

\begin{corollary}[rank-one per-sample noise]\label{cor:rank1}
If $\zeta$ has rank one almost surely (per-sample gradients $\delta x^\top$ of a linear layer), covariance $\sigma^2I$
and satisfies Assumption~\ref{ass:N} (Cram\'er's condition is not automatic for a law carried by the rank-one variety),
then $\zeta\zeta^\top\zeta=\nF\zeta^2\zeta$ and
$-\Med_\ast(\xi_b)=\frac{3+2q_{K,d}}{(K+2)(d+2)}\cdot\frac{\E[\zeta\nF\zeta^2]}{2\sigma^2b}+o(\frac1b)$. For the same noise a
Frobenius-normalized step has the offset $\frac{1}{Kd+2}\cdot\frac{\E[\zeta\nF\zeta^2]}{2\sigma^2b}$: orthogonalization
multiplies the first-order offset of normalization by $\frac{(3+2q_{K,d})(Kd+2)}{(K+2)(d+2)}\in(1,r)$ (it exceeds
$\frac{3(Kd+2)}{(K+2)(d+2)}>1$ and is below $\frac{Kd+2}{\max(K,d)+2}<r$): $1.37$ for $2\times2$ matrices and about
$2\kappa(\gamma)\,r$ for large ones ($0.8\,r$ for square matrices).
\end{corollary}

\begin{corollary}[what is seen, what is annihilated; no domination]\label{cor:compare}
(i) To first order a polar step sees the third cumulant only through $a_1\cA+a_2\cB$ ($Kd$ numbers), the normalized step
only through $\cA$, sign descent only through the $Kd$ own third cumulants $K_3^{(ij),(ij),(ij)}$. (ii) \emph{Sign-blind
noise biases polar steps:} if $\zeta$ has Gaussian entries except $\zeta_{11}$, which has zero third cumulant but
$\E[\zeta_{11}\zeta_{12}^2]=\tau\ne0$, sign descent has no first-order offset and
$-\Med_\ast(\xi_b)=\frac{(a_1+a_2)\tau}{6\sigma^2b}e_1e_1^\top+o(1/b)$. (iii) \emph{Polar-blind noise biases sign descent:}
every $K_3$ with $a_1\cA[K_3]+a_2\cB[K_3]=0$ and a nonzero own cumulant, e.g.\ for $K,d\ge2$ and $a_1\ne0$ a tensor supported on
$\{(11),(22)\}$ with $K_3^{(11)^3}=1$, $K_3^{(11),(22),(22)}=-\frac{a_1+a_2}{a_1}$ and $K_3^{(22)^3}=K_3^{(22),(11),(11)}=0$.
(iv) \emph{One skewed factor:} if $\zeta=\epsilon M+\zeta'$ with $\epsilon$ a scalar of third cumulant $\kappa_3(\epsilon)$,
$M$ fixed, $\zeta'$ independent and symmetric, and total covariance $\sigma^2I$, then
$g_\ast=\frac{\kappa_3(\epsilon)}{6\sigma^2}(a_1\nF M^2M+a_2MM^\top M)$ and $\inner{g_\ast}{\kappa_3(\epsilon)M}>0$: the limit
gradient is positively aligned with the skewed direction, as in the scalar case.
\end{corollary}
\begin{proof}
(i) is Theorems~\ref{thm:cf} and~\ref{thm:two} and Remark~\ref{rem:cf}. (ii) The only nonzero entries of $K_3$ are the
permutations of $((11),(12),(12))$, whence $\cA=\cB=\tau e_1e_1^\top$. Such a law exists:
$\zeta_{11}=\rho(\zeta_{12}^2-1)/\sqrt2+\sqrt{1-\rho^2}\,\epsilon$ with $\epsilon$ an independent centred Gamma variable
whose skewness cancels that of the first term (the law of \texttt{c03}, with a density). (iii) For this tensor
$\cA_{11}=1+K_3^{(11),(22),(22)}$, $\cB_{11}=1$ and all other entries of $\cA,\cB$ vanish (the pattern
$(kj),(kl),(il)$ of $\cB$ is never a permutation of $((11),(22),(22))$), so
$a_1\cA+a_2\cB=(a_1+a_2+a_1K_3^{(11),(22),(22)})e_1e_1^\top=0$; a small multiple of this $K_3$ is the third cumulant of a law with
covariance $I$ (the pairs of covariance and third moment of laws with a density near the Gaussian fill a neighbourhood
of $(I,0)$). (iv) Cumulants add over independent summands and
$a_1\nF M^4+a_2\tr(MM^\top)^2>0$ by Theorem~\ref{thm:two}(d).
\end{proof}

\begin{remark}[sign of $a_1$]\label{rem:a1}
$a_1<0$ is equivalent to $q_{K,d}>\frac{(K-1)(d-1)}{2(K+d+1)}$. It holds on the fifteen sizes computed (nine are shown in
Table~\ref{tab:const}), in the proportional limit (numerically $\kappa(\gamma)>\frac1{2(1+\gamma)}$ on $(0,1]$, with equality
at $0^+$), and for fixed $K$ as $d\to\infty$, where a second-order expansion of the singular values around $\sqrt d$ gives
$q=\frac{K-1}2-\frac{(K-1)(K+2)}{16d}+o(\frac1d)$ (heuristic; $0.4961$ against the computed $0.4960$ for $2\times64$),
against the threshold
$\frac{K-1}2-\frac{(K-1)(K+2)}{2d}+o(\frac1d)$; we have no proof for all sizes. In \eqref{eq:two} the two moments therefore
pull in opposite directions: writing $\zeta=U\operatorname{diag}(s)V^\top$,
$a_1\zeta\nF\zeta^2+a_2\zeta\zeta^\top\zeta=U\operatorname{diag}\big(s_i(a_2s_i^2-|a_1|\,|s|^2)\big)V^\top$, so a singular
component of a sample contributes with the sign of its skew only if it carries more than the fraction $|a_1|/a_2$ of the
energy of that sample.
\end{remark}

\section{Parameter-dependent noise}\label{sec:state}

Let the stochastic gradient be $\hat g_b(W)=\frac1b\sum_{i\le b}g(W;\omega_i)$, $\nabla F(W)=\E g(W;\omega)$, and let
$\Sigma(W)$, $K_3(W)$ be the covariance and third cumulant of $g(W;\omega)-\nabla F(W)$, $X_b^W$ the normalized sum. The
mean field is $h_b(W)=\E\,o(\hat g_b(W))$; with parameter-dependent noise it is not a gradient field and
\cite[Theorem~2]{manuscript} does not apply.

\begin{assumption}[local regularity and uniform expansion]\label{ass:U}
On a bounded neighbourhood $\mathcal U$ of $W_\ast$: (i) $F$ is $C^1$, differentiable twice at $W_\ast$,
$\nabla F(W_\ast)=0$, $H=\nabla^2F(W_\ast)$ is invertible and $|\nabla F(W)|\ge c|W-W_\ast|$; (ii) $\Sigma(W)$ and
$K_3(W)$ are continuous at $W_\ast$, bounded, and $\Sigma(W)\succeq\lambda I$, $\lambda>0$; (iii) the conclusion of
Lemma~\ref{lem:edge} holds uniformly:
$\sup_{W\in\mathcal U}\sup_h|\E\,o(X_b^W+h)-\Phi_{\Sigma(W)}(h)-\frac1{6\sqrt b}D^3\Phi_{\Sigma(W)}(h)[K_3(W)]|=o(b^{-1/2})$.
\end{assumption}

\begin{theorem}[localization of the zeros]\label{thm:state}
Under Assumptions~\ref{ass:O} and~\ref{ass:U}, let $g_\ast=-\frac16J_{\Sigma_\ast}^{-1}T_{\Sigma_\ast}[K_{3\ast}]$ be
\eqref{eq:gstar} at $W_\ast$. For every $\varepsilon>0$ there are $b_0$ and $\delta>0$ such that for $b\ge b_0$ every
$W\in\mathcal U$ with $|h_b(W)|\le\delta b^{-1/2}$, in particular every zero of the mean field in $\mathcal U$, satisfies
\[
\Big|W-W_\ast-\frac1b\,H^{-1}g_\ast\Big|\le\frac\varepsilon b .
\]
If moreover $h_b$ is continuous on $\mathcal U$, it has a zero in $\mathcal U$ for $b$ large.
\end{theorem}
\begin{proof}
By homogeneity and (iii), uniformly in $W\in\mathcal U$,
\begin{equation}\label{eq:hb}
h_b(W)=\Phi_{\Sigma(W)}\big(\sqrt b\,\nabla F(W)\big)+\tfrac1{6\sqrt b}D^3\Phi_{\Sigma(W)}\big(\sqrt b\,\nabla F(W)\big)[K_3(W)]+o(b^{-1/2}).
\end{equation}
\emph{Step 1.} By Lemma~\ref{lem:gauss}, uniformly over the compact set of covariances,
$|\Phi_{\Sigma(W)}(h)|\ge c_0\min(|h|,1)$ and the second term of \eqref{eq:hb} is at most $C_1b^{-1/2}$. Hence
$|h_b(W)|\le\delta b^{-1/2}$ forces $c_0\min(\sqrt b|\nabla F(W)|,1)\le(C_1+\delta+o(1))b^{-1/2}$, i.e.\
$|\nabla F(W)|\le C_2/b$ and $|W-W_\ast|\le C_3/b$ for $b$ large. \emph{Step 2.} Write $W=W_\ast+w/b$, $|w|\le C_3$. Then
$\sqrt b\,\nabla F(W)=b^{-1/2}(Hw+o(1))$, $\Sigma(W)\to\Sigma_\ast$, $K_3(W)\to K_{3\ast}$, and $J_\Sigma$, $T_\Sigma$ are
continuous in $\Sigma$, so \eqref{eq:hb} reads $h_b(W)=b^{-1/2}\big(J_{\Sigma_\ast}Hw+\frac16T_{\Sigma_\ast}[K_{3\ast}]+o(1)\big)$,
uniformly in $|w|\le C_3$. If $|h_b(W)|\le\delta b^{-1/2}$ then $|w-H^{-1}g_\ast|\le\|(J_{\Sigma_\ast}H)^{-1}\|(\delta+o(1))$.
\emph{Existence.} The continuous field $w\mapsto\sqrt b\,h_b(W_\ast+w/b)$ converges uniformly on a ball around
$H^{-1}g_\ast$ to the affine isomorphism $w\mapsto J_{\Sigma_\ast}H(w-H^{-1}g_\ast)$, which has degree $\pm1$ there.
\end{proof}

\begin{proposition}[when Assumption~\ref{ass:U}(iii) holds]\label{prop:lin}
Suppose $g(W;\omega)-\nabla F(W)=L_W\chi(\omega)$ on $\mathcal U$, where $\chi\in\R^p$ satisfies Assumption~\ref{ass:N}
and $W\mapsto L_W$ is continuous with $L_W:\R^p\to E$ onto. Then (ii) and (iii) hold on a neighbourhood of $W_\ast$.
This covers least squares, $g=(Wx-y)x^\top=(W-W_0)(xx^\top)-\epsilon x^\top$, with
$\chi=(xx^\top-\E xx^\top,\ \epsilon x^\top)$.
\end{proposition}
\begin{proof}
Apply Theorem~20.1 of \cite{br} to $\chi$ and to $f_{W,h}(x)=o(L_Wx+h)$: its oscillation on a ball of radius
$\varepsilon$ is that of $o$ on a ball of radius $\|L_W\|\varepsilon$ around $L_Wx+h$, and $L_WZ_\chi\sim
N(0,\Sigma(W))$ with $\Sigma(W)=L_W\Sigma_\chi L_W^\top$ in a compact set of positive definite matrices, so
Assumption~\ref{ass:O} bounds the oscillation term uniformly; the remainder $\delta_b$ depends on $\chi$ only. The
Gaussian and first Edgeworth terms of $\chi$ push forward to those of covariance $\Sigma(W)$ and cumulant
$L_W^{\otimes3}K_3^\chi=K_3(W)$: with $u(s)=\E f_{W,h}(Z_\chi+s)=\Phi_{\Sigma(W)}(h+L_Ws)$,
$-\int f_{W,h}K_3^\chi\cdot\partial^3\varphi_{\Sigma_\chi}=D^3u(0)[K_3^\chi]=D^3\Phi_{\Sigma(W)}(h)[L_W^{\otimes3}K_3^\chi]$.
\end{proof}

\begin{remark}
(i) The Hessian enters only through $H^{-1}g_\ast$: at first order the variation of the noise law with $W$ does not move
the zero (numerically the relative difference between the zeros with the true and with the frozen noise law decreases
from $6\%$ at $b=16$ to $0.6\%$ at $b=256$, Table~\ref{tab:checks}).
(ii) Uniform one-term expansions over families of laws under uniform moment and weak Cram\'er conditions are in
\cite{song}; they would give (iii) for smooth models, a verification we have not carried out. (iii) Cram\'er's condition
for $\chi$ holds when $(x,\epsilon)$ has a density, by the classical argument for analytic statistics \cite{bg}; this
citation is from memory.
\end{remark}

\section{Momentum}\label{sec:mom}

Consider $M_t=\beta M_{t-1}+(1-\beta)(\nabla F(W_t)+\xi_t)$, $W_{t+1}=W_t-\eta_t\polar(M_t)$, with $\beta\in[0,1)$ fixed,
$\xi_t$ i.i.d.\ copies of $\xi$, and let $N_\beta=(1-\beta)\sum_{k\ge0}\beta^k\xi_k$ be the stationary filtered noise.

\begin{assumption}\label{ass:M}
$0<\muop\le\Lop<\infty$ (strong convexity and smoothness of $F$ for the spectral norm);
$\E\nF\xi^2<\infty$; there is $C$ with $\sup_{Y\in E}\Prob(\sigma_r(Y+\xi)\le t)\le Ct$ for all $t>0$ (true for a Gaussian
$\xi$ by Lemma~\ref{lem:osc}, and for densities whose conditional row densities are bounded); and
$\eta_t=\eta_0(t+t_0)^{-\theta}$ with $\theta\in(\frac12,1]$.
\end{assumption}

\begin{theorem}[convergence with momentum]\label{thm:momconv}
Under Assumption~\ref{ass:M}, $\nabla F(W_t)$ converges almost surely to a point of
$\arg\min\psi_{N_\beta}=-\Med_\ast(N_\beta)$. In particular $W_t\to W_\ast$ if $\Med_\ast(N_\beta)=\{0\}$ (e.g.\ $\xi$
symmetric with $0$ in its support), and otherwise the limit is offset by a nuclear median of the \emph{filtered} noise.
\end{theorem}
\begin{proof}
Write $\cF_t=\sigma(\xi_1,\dots,\xi_t)$, $G_t=\nabla F(W_t)\in\cF_{t-1}$, $O_t=\polar(M_t)$, $\psi=\psi_{N_\beta}$,
$\psi_\ast=\min\psi$, fix $G^\st\in\arg\min\psi$ and $V_t=D_{F^\ast}(G^\st,G_t)=D_F(W_t,W^\st)$. As in
\cite[Theorem~2]{manuscript}, the three-point identity gives, pathwise,
\begin{equation}\label{eq:3pt}
V_{t+1}=V_t-\eta_tX_t+D_F(W_{t+1},W_t),\qquad X_t=\inner{O_t}{G_t-G^\st},\qquad0\le D_F(W_{t+1},W_t)\le\tfrac{\Lop}2\eta_t^2,
\end{equation}
and $\nnuc{G-G^\st}^2\le2\Lop D_{F^\ast}(G^\st,G)$, $\nnuc{G_{t}-G_{s}}\le\Lop\sum_{s\le u<t}\eta_u$.

\emph{Decomposition.} $M_t=\bar G_t+N_t$ with $N_t=(1-\beta)\sum_{k<t}\beta^k\xi_{t-k}$ and $\bar G_t$ the same average
of the gradients (plus $\beta^tM_0$); $\nnuc{\bar G_t-G_t}\le\varepsilon_t$ with a deterministic
$\varepsilon_t\le C\eta_t$. Let $m=m_t=\lceil c_\beta\log t\rceil$ with $\beta^{m_t}\le t^{-2}$, and split
$N_t=N_t^{(m)}+\beta^mN_{t-m}$, where $N_t^{(m)}=(1-\beta)\sum_{k<m}\beta^k\xi_{t-k}$ is independent of $\cF_{t-m}$. Put
$Y_t=G_{t-m}+N_t^{(m)}$, $\tilde O_t=\polar(Y_t)$, $\tilde X_t=\inner{\tilde O_t}{G_{t-m}-G^\st}$. Then
$M_t=Y_t+\Delta_t$ with $\nF{\Delta_t}\le\rho_t:=r_t+\beta^m\nF{N_{t-m}}$, $r_t=\Lop m_t\eta_{t-m_t}+\varepsilon_t$
deterministic, and $\rho_t$ is $\cF_{t-m}$-measurable.

\emph{The lagged drift.} Conditionally on $\cF_{t-m}$, $N_t^{(m)}$ has the law of the truncated filter, so
$c_t:=\E[\tilde X_t\mid\cF_{t-m}]=\inner{\Phi_{N^{(m)}}(G_{t-m})}{G_{t-m}-G^\st}\ge\psi_{N^{(m)}}(G_{t-m})-\psi_{N^{(m)}}(G^\st)
\ge p_t-2\beta^m\E\nnuc{N_\beta}$, where $p_t=\psi(G_{t-m})-\psi_\ast\ge0$ (the differences $\psi(G)-\psi(G^\st)$ computed with the two noises differ by at most
$2\beta^m\E\nnuc{N_\beta}$).

\emph{The perturbation.} By \cite{lisun}, $\nF{O_t-\tilde O_t}\le2\sqrt r\min(1,2\rho_t/\sigma_r(Y_t))$. Conditionally on
$\cF_{t-m}$, $Y_t=(\text{independent part})+(1-\beta)\xi_t$, so $\Prob(\sigma_r(Y_t)\le u\mid\cF_{t-m})\le Cu/(1-\beta)$
and $\E[\min(1,2\rho_t/\sigma_r(Y_t))\mid\cF_{t-m}]\le\int_0^1\min(1,C'\rho_t/u)du\le C_\kappa\rho_t^\kappa$ for any fixed
$\kappa\in(0,1)$; choose $\kappa$ with $\theta(1+\kappa)>1$. Hence
\[
\E\big[|X_t-\tilde X_t|\ \big|\ \cF_{t-m}\big]\le C\rho_t^\kappa\,\nnuc{G_{t-m}-G^\st}+\Lop m_t\eta_{t-m_t}
\le C\big(r_t^\kappa+\beta^{m\kappa}\nF{N_{t-m}}^\kappa\big)\sqrt{2\Lop V_{t-m}}+\Lop m_t\eta_{t-m_t}.
\]
\emph{Moment bound.} From \eqref{eq:3pt}, $|V_{t+1}-V_t|\le C\eta_t(1+V_t)$, so $1+V_{t-m}\le2(1+V_t)$ for $t$ large. Taking
expectations in \eqref{eq:3pt}, with $v_t=\E V_t$, the Cauchy--Schwarz inequality and $\E\nF{N}^{2\kappa}<\infty$,
$v_{t+1}\le v_t+C\eta_t\delta_t(1+v_t)-\eta_t\E p_t$, $\delta_t=r_t^\kappa+\beta^{m_t\kappa}+m_t\eta_{t-m_t}+\eta_t$, and
$\sum_t\eta_t\delta_t<\infty$ because $\eta_tr_t^\kappa=O(t^{-\theta(1+\kappa)}\log t)$. Hence $\sup_tv_t<\infty$.

\emph{Three convergent series.} (1) $\sum_t\eta_t|X_t-\tilde X_t|<\infty$ almost surely, its expectation being finite.
(2) $\sum_t\eta_t(\tilde X_t-c_t)$ converges almost surely: write
$\tilde X_t-c_t=\sum_{j<m_t}D_{t,j}$, $D_{t,j}=\E[\tilde X_t|\cF_{t-j}]-\E[\tilde X_t|\cF_{t-j-1}]$; for fixed $j$,
$\sum_t\eta_tD_{t,j}\mathbf 1_{m_t>j}$ is a martingale in $t$ for the filtration $(\cF_{t-j})_t$ with
$\E D_{t,j}^2\le4\E\tilde X_t^2\le8\Lop v_{t-m}\le C$, so by Doob's inequality and Minkowski's inequality
$\|\sup_{T\ge T_0}|\sum_{T_0<t\le T}\eta_t(\tilde X_t-c_t)|\|_2\le C\sum_{j\ge0}\big(\sum_{t>\max(T_0,e^{j/c_\beta})}\eta_t^2\big)^{1/2}$,
which tends to $0$ as $T_0\to\infty$ since $\theta>\frac12$. (3) $\sum_tD_F(W_{t+1},W_t)<\infty$. Summing
\eqref{eq:3pt}: $V_{T+1}+\sum_{t\le T}\eta_tc_t$ converges almost surely; as $V\ge0$ and
$\eta_tc_t\ge\eta_tp_t-2\eta_t\beta^{m_t}\E\nnuc{N_\beta}$ with a summable negative part, both $\sum_t\eta_tp_t<\infty$ and
the convergence of $V_t$ follow.

\emph{Conclusion.} $\sum_t\eta_t=\infty$ gives $\liminf_tp_t=0$; $V_t$ bounded gives a cluster point $\hat G$ of
$G_{t-m_t}$, hence of $G_t$ (since $\nnuc{G_t-G_{t-m_t}}\to0$), with $\psi(\hat G)=\psi_\ast$. Running the argument for a
countable dense subset of the compact set $\arg\min\psi$ and using the continuity of $G^\st\mapsto D_{F^\ast}(G^\st,G)$,
$D_{F^\ast}(\hat G,G_t)$ converges, and its limit is $0$ along a subsequence: $G_t\to\hat G$.
\end{proof}

\begin{theorem}[first-order bias with a linear filter]\label{thm:mombias}
Let $(w_k)_{k\ge0}$ be deterministic weights with $w_0\ne0$, $\sum_k|w_k|<\infty$, $\rho_p=\sum_kw_k^p$, and
$N_w=\sum_kw_k\xi_b^{(k)}$ with independent copies $\xi_b^{(k)}$ of $\xi_b$. Under Assumptions~\ref{ass:N}
and~\ref{ass:O} with $\E|\zeta|^4<\infty$, for $g_\ast$ of \eqref{eq:gstar},
\[
-\Med_\nu(N_w)=\frac{\rho_3}{\rho_2}\cdot\frac{g_\ast}b+o\Big(\frac1b\Big)\qquad(b\to\infty):
\]
the bias is that of the unfiltered noise at the batch size $b_{\mathrm{eff}}=b\,\rho_2/\rho_3$. For the exponential
moving average ($w_k=(1-\beta)\beta^k$), $b_{\mathrm{eff}}=b\,\frac{1+\beta+\beta^2}{1-\beta^2}$; for the Nesterov
look-ahead of Muon ($w_0=1-\beta^2$, $w_k=(1-\beta)\beta^{k+1}$),
$\rho_p=(1-\beta^2)^p+(1-\beta)^p\beta^{2p}/(1-\beta^p)$. The covariance of $N_w$ is that of batch size $b/\rho_2$
($b\frac{1+\beta}{1-\beta}$ for the average).
\end{theorem}
\begin{proof}
Let $X=\sqrt b\,N_w=w_0X_0+R$, $R=\sum_{k\ge1}w_kX_k$, $X_k=\sqrt b\,\xi_b^{(k)}$ independent. Conditionally on $R$,
Lemma~\ref{lem:edge} applied to $X_0$ and the homogeneity of $o$ give, uniformly in $h$,
$\E\,o(X+h)=\E\,\Phi_\Sigma\big(\tfrac{R+h}{w_0}\big)+\frac1{6\sqrt b}\E\,D^3\Phi_\Sigma\big(\tfrac{R+h}{w_0}\big)[K_3]+o(b^{-1/2})$
(for $w_0<0$ use the oddness of $o$). $R=\sum_{k\ge1}\sum_{i\le b}\frac{w_k}{\sqrt b}\zeta_{k,i}$ is a sum of independent
centred terms with coefficients $a_{k,i}=w_k/\sqrt b$. For a function $\Psi$ with bounded derivatives up to order six,
Lindeberg's replacement with a third-order Taylor expansion gives
\[
\Big|\E\Psi(R)-\E\Psi(Z_R)-\tfrac16\textstyle\sum a_{k,i}^3\,\E D^3\Psi(Z_R)[K_3]\Big|\le
C_\Psi\Big(\textstyle\sum a_{k,i}^4+\big(\sum|a_{k,i}|^3\big)^2\Big)=O(b^{-1}),
\]
with $Z_R\sim N(0,(\rho_2-w_0^2)\Sigma)$: replacing $\zeta_{k,i}$ by a Gaussian changes the expectation by
$\frac{a^3}6\E D^3\Psi(U)[K_3]+O(a^4)$, $U$ the partially replaced sum, and $\E D^3\Psi(U)=\E D^3\Psi(Z_R)+O(\sum|a|^3+a^2)$ by
the same argument at one order less. Apply this to $\Psi=\Phi_\Sigma((\cdot+h)/w_0)$ and, at order zero, to
$D^3\Phi_\Sigma((\cdot+h)/w_0)[K_3]$; since $\E\,\Phi_\Sigma((Z_R+h)/w_0)=\E\,o(w_0Z_0+Z_R+h)=\Phi_{\rho_2\Sigma}(h)$ and
$D^3$ of this function of $h$ collects the factor $w_0^3$,
\[
\sup_h\Big|\E\,o(X+h)-\Phi_{\rho_2\Sigma}(h)-\frac{\rho_3}{6\sqrt b}D^3\Phi_{\rho_2\Sigma}(h)[K_3]\Big|=o(b^{-1/2}).
\]
The proof of Theorem~\ref{thm:cf} then applies with $J_{\rho_2\Sigma}=\rho_2^{-1/2}J_\Sigma$,
$T_{\rho_2\Sigma}=\rho_2^{-3/2}T_\Sigma$ and $\rho_3K_3$ in the place of $K_3$.
\end{proof}

\begin{remark}\label{rem:mom}
(i) Theorems~\ref{thm:momconv} and~\ref{thm:mombias} together: Muon-type steps with a fixed momentum and Robbins--Monro
lengths converge to a point whose gradient is $\frac{1-\beta^2}{1+\beta+\beta^2}\frac{g_\ast}b+o(1/b)$. As
$\beta\to1$ the factor is $\frac23(1-\beta)$, against $\frac12(1-\beta)$ for the variance: the skewness-driven bias is
governed by another effective batch size than the one published for the variance and the curvature-driven displacement
\cite{paquette,chen}. The cumulant ratios of the filtered variable appear in \cite[Prop.~6.11]{ddj} for another
functional (a Lipschitz Newton--Schulz map in dimension one, where the offset has another order), and the scalar factor
with a frozen iterate is in the archived notes of the project; the theorem for the coupled dynamics, the exact polar
factor and matrices is new. (ii) The regime $\beta\to1$ at fixed $b$ (even $b=1$) is not covered by
Theorem~\ref{thm:mombias}; numerically, at $b=1$ the formula with $b_{\mathrm{eff}}$ is within $8\%$, $3\%$, $1\%$ for the
average at $\beta=0.5,0.9,0.95$ and within $10\%$, $7\%$ for the Nesterov weights at $\beta=0.9,0.95$
(Table~\ref{tab:checks}). A proof needs an Edgeworth expansion for
geometrically weighted sums; Theorem~20.6 of \cite{br} on non-identically distributed summands is the natural tool.
\end{remark}

\section{Sampling without replacement}\label{sec:wor}

\begin{proposition}[complement duality]\label{prop:wor}
Let $v_1,\dots,v_n\in E$ with $\sum_iv_i=0$ and let $\xi_b$ be the mean of a uniformly random subset of size $b$. Then
$\xi_{n-b}$ has the law of $-\frac b{n-b}\xi_b$, hence for every geometry
\[
\Med_\nu(\xi_{n-b})=-\frac b{n-b}\,\Med_\nu(\xi_b),\qquad 0\in\Med_\nu(\xi_{n/2})\ (n\text{ even}).
\]
With the finite-population cumulants, $\mathrm{Cov}(\xi_b)=\frac{n-b}{b(n-1)}\Sigma$ and
$K_3(\xi_b)=\frac{(n-b)(n-2b)}{b^2(n-1)(n-2)}K_3$, the first-order law reads \emph{(formal)}
$-\Med_\nu(\xi_b)\approx\frac{n-2b}{n-2}\cdot\frac{g_\ast}b$, which is consistent with the duality and changes sign at $b=n/2$.
\end{proposition}
\begin{proof}
If $S$ is uniform among subsets of size $b$, its complement is uniform among subsets of size $n-b$ and
$\sum_{i\notin S}v_i=-\sum_{i\in S}v_i$. The median is odd and positively homogeneous. For $b=n/2$ the law of $\xi_b$ is
symmetric and $\psi_{\xi_b}$ is even and convex.
\end{proof}

\section{Numerical checks}\label{sec:num}

All scripts are in \texttt{experiments/bias\_law/checks/}, logs in \texttt{results/bias\_law/checks/} (\texttt{c01.log}, \dots, \texttt{c07.log};
an independent re-derivation of the algebra of Theorem~\ref{thm:two} and of its constants is in
\texttt{checks/review\_verify.py}). Zeros of mean fields are computed by root
finding on $S$ sampled batch means with common random numbers ($S$ between $6\cdot10^5$ and $4\cdot10^6$); $J_\Sigma$,
$T_\Sigma$ of the general law by Monte Carlo ($4\cdot10^6$ antithetic Gaussian pairs).

\begin{table}[h]
\centering\small
\begin{tabular}{@{}rrrrrrrr@{}}
\toprule
$K$ & $d$ & $q_{K,d}$ & $\frac{r-1}2$ & $a_1$ & $a_2$ & $a_1+a_2$ & $\max(K,d)(a_1+a_2)$\\
\midrule
2 & 2 & 0.3204 & 0.5 & $-0.4133$ & 1.0959 & 0.6827 & 1.365\\
2 & 3 & 0.3935 & 0.5 & $-0.2041$ & 0.7722 & 0.5680 & 1.704\\
3 & 5 & 0.8425 & 1.0 & $-0.0768$ & 0.4784 & 0.4016 & 2.008\\
4 & 4 & 1.0680 & 1.5 & $-0.0947$ & 0.5227 & 0.4280 & 1.712\\
8 & 8 & 2.6353 & 3.5 & $-0.0249$ & 0.2730 & 0.2481 & 1.985\\
16 & 16 & 5.8149 & 7.5 & $-0.0065$ & 0.1420 & 0.1355 & 2.167\\
64 & 64 & 24.995 & 31.5 & $-0.0004$ & 0.0369 & 0.0365 & 2.336\\
8 & 32 & 3.3508 & 3.5 & $-0.0024$ & 0.0880 & 0.0856 & 2.739\\
32 & 128 & 14.940 & 15.5 & $-0.0002$ & 0.0225 & 0.0223 & 2.857\\
\bottomrule
\end{tabular}
\caption{Constants of Theorem~\ref{thm:two} (Monte Carlo for $q$; \texttt{c02}). The column $a_1+a_2$ reproduces the
Monte-Carlo table of $a_{K,d}/c_{K,d}$ of the archived notes ($0.683$, $0.568$, $0.402$, $0.428$, $0.250$, $0.136$,
$0.086$). Limits: $\theta(1)=2.4$, $\theta(\frac14)=2.898$. The Stein identity of (b) holds within Monte-Carlo
error.}\label{tab:const}
\end{table}

\begin{table}[h]
\centering\small
\begin{tabular}{@{}p{0.31\textwidth}p{0.64\textwidth}@{}}
\toprule
Check (script) & Outcome\\
\midrule
Two-invariant structure (\texttt{c01}) & Monte-Carlo tensor $\E[\polar(Z)\otimes\He_3(Z)]$ for $2\times2$, $2\times3$:
residual after projection on $\{\cA,\cB\}$ $0.5\%$ and $1\%$ of its norm (Monte-Carlo level).\\
I.i.d.\ $\mathrm{Exp}-1$ entries, $2\times3$ (\texttt{c03}) & $g_\ast=0.189\,\ones\ones^\top$; $b\,G_b$ within $1.5\%$, $2.1\%$,
$4.9\%$ for $b=4,16,64$ (Monte-Carlo error $1.5$--$8\%$).\\
Sign-blind co-skewness, Cor.~\ref{cor:compare}(ii) (\texttt{c03}) & predicted $0.080\,e_1e_1^\top$; measured $(b\,G_b)_{11}=0.068$,
$0.077$, $0.085$ at $b=4,16,64$; sign descent: $|b\cdot\text{median}|\le0.02$, compatible with $0$.\\
One skewed factor of full rank (\texttt{c03}) & $g_\ast$ uses $a_1$ and $a_2$ separately; relative error $10$--$14\%$ at the
Monte-Carlo level; cosine $0.990$--$0.998$.\\
Rank-one noise, isotropic (\texttt{c03}) & convergence is slow: $b\,G_b/g_\ast\approx-0.02$, $0.32$, $0.68$, $0.87$, $0.96$ for
$b=4,16,64,256,1024$ (the offset has the opposite sign at $b=4$).\\
Rank-one, Kronecker covariance, general law (\texttt{c03}) & relative error $0.62$, $0.26$, $0.05$ at $b=16,64,256$; cosine
$\ge0.999$ from $b=16$.\\
Plateau of the manuscript, Fig.~3(c) (\texttt{c07}) & finite-sum least squares, $n=2000$: predicted
$\nnuc{g_\ast}=0.3139\pm0.0005$; the manuscript measures $b\nnuc{\nabla F(W_b)}=0.276$, $0.285$, $0.303$, $0.308$ at
$b=4,8,16,32$.\\
Population least squares, state-dependent (\texttt{c06}) & $b(W_b-W_0)$ against $H^{-1}g_\ast$: relative error $0.18$,
$0.10$, $0.12$ at $b=16,64,256$ (Monte-Carlo error $0.02$, $0.04$, $0.15$); the zeros with the true and with the frozen
noise law differ by $6\%$, $1.6\%$, $0.6\%$.\\
Momentum, frozen iterate (\texttt{c04}) & measured$/$first-order: $0.92$ ($\beta=0.5,b=1$), $1.00$ ($0.5,8$), $0.97$
($0.9,1$), $1.02$ ($0.9,4$), $1.01$ ($0.95,1$); Nesterov: $0.90$ ($0.9,1$), $0.93$ ($0.95,1$). The variance-based
effective batch would give $0.75$--$0.78$ of the first-order value for the average.\\
Momentum, coupled dynamics (\texttt{c04}) & tail-averaged gradient of $2000$ runs within $3.3\%$, $2.1\%$ ($\beta=0.9$, two
step scales) and $0.5\%$ ($\beta=0.5$) of the zero of the frozen stationary mean field.\\
Without replacement (\texttt{c05}) & duality exact to $10^{-9}$ on enumerated populations; $n=400$: coefficient of $G_b$ on
$g_\ast$ equals $\frac{n-2b}{b(n-2)}$ within $2$--$12\%$ for $b\le150$ (at $b=8$ the measured $0.1243$ does not separate
it from $1/b$), vanishes at $b=200$ ($10^{-5}$) and is negative
beyond ($-0.0011$, $-0.0016$, $-0.0022$ at $b=250,300,350$; predicted $-0.0010$, $-0.0017$, $-0.0022$).\\
\bottomrule
\end{tabular}
\caption{Numerical checks.}\label{tab:checks}
\end{table}

\section{Novelty ledger, limits, open problems}\label{sec:ledger}

\paragraph{Closest results} (details and the scope of the search in \texttt{theory/bias\_law/notes/novelty\_ledger.md}).
\emph{Manuscript and repository:} Theorem~2 and Propositions~3--4 of \cite{manuscript}; in the archived notes, the formal
operator formula $-\frac1b\mathcal J^{-1}\mathcal T$, the Monte-Carlo table of $a/c$, $a/c=3/(d+2)$ for $K=1$ and the
scalar momentum factors with a frozen iterate. \emph{Literature:} the linear response $c_{K,d}$ of Gaussian-smoothed
polar maps \cite{xie,chen} (no third-order tensor, no non-Gaussian input); the geometric median of the law of a mean
\cite{minsker} and its expansion in the dimension \cite{schwank}; Muon in dimension one with momentum, where every
limit is a zero of the mean field of the filtered noise, for Lipschitz Newton--Schulz maps and assuming convergence
\cite{ddj}; variance-based effective batch sizes \cite{paquette,chen}. Not found: Theorem~\ref{thm:cf} at fixed
dimension for a general norm, Theorem~\ref{thm:two} and its corollaries, Theorem~\ref{thm:state},
Theorem~\ref{thm:momconv}, the effective batch $b\rho_2/\rho_3$ as a theorem, Proposition~\ref{prop:wor}. A failure to
find is not a proof of absence: 23 queries and about thirty sources were opened; no citation-graph search was made.

\paragraph{What is not proved.} $a_1<0$ for all sizes (Remark~\ref{rem:a1}); the finite-sum (discrete) case of
Theorem~\ref{thm:cf} and the formal part of Proposition~\ref{prop:wor}; Assumption~\ref{ass:U}(iii) beyond linear
models; the regime $\beta\to1$ at fixed $b$; a closed form of $q_{K,d}$ (it is a one-point integral of the real Laguerre
ensemble; $q_{2,2}$ and $q_{2,d}$ are elementary). No closed form is available for anisotropic covariances, where
$J_\Sigma,T_\Sigma$ are Gaussian integrals evaluated by Monte Carlo; for the rank-one, heavy-tailed laws tested the
first-order regime starts only at $b$ of the order of $10^2$. Constant learning rates, Newton--Schulz maps (for which
the offset has another order \cite{ddj}) and non-convex objectives are outside the note.

\paragraph{Open problems.} (1) Prove $a_1<0$, i.e.\ $q_{K,d}>\frac{(K-1)(d-1)}{2(K+d+1)}$. (2) A Kronecker covariance
$\Sigma_L\otimes\Sigma_R$ breaks the symmetry used in Theorem~\ref{thm:two}; is there a deterministic equivalent of
$J_\Sigma^{-1}T_\Sigma$ in the proportional limit? (3) The next term: for symmetric noise the offset vanishes, and for
$K_3=0$ it is $O(b^{-2})$; which cumulants drive it? (4) The sign reversal of the offset at small batch for rank-one
noise (Table~\ref{tab:checks}) is not explained by a one-term expansion.

\begin{thebibliography}{99}\small
\bibitem{manuscript} Project ``Accelerating Variance Reduction'', \emph{From SAGA to Muon-type steps: no norm repairs the
logistic worst case, and variance reduction acts on the fixed point}, \texttt{paper/paper\_saga\_to\_muon.tex}, and the
archived notes \texttt{archive/muon/notes\_spectral\_theory.tex} (Proposition ``size of the offset'').
\bibitem{br} R.~N. Bhattacharya and R.~R. Rao, \emph{Normal Approximation and Asymptotic Expansions}, Wiley 1976; SIAM
Classics 2010. Theorems 20.1 and 20.6.
\bibitem{angstpoly} J. Angst and G. Poly, A weak Cram\'er condition and application to Edgeworth expansions,
\emph{Electron.\ J. Probab.} 22 (2017), arXiv:1511.08750.
\bibitem{song} K. Song, A uniform-in-$P$ Edgeworth expansion under weak Cram\'er conditions, arXiv:1806.01431.
\bibitem{bg} R.~N. Bhattacharya and J.~K. Ghosh, On the validity of the formal Edgeworth expansion, \emph{Ann.\ Statist.}
6 (1978).
\bibitem{lisun} W. Li and W. Sun, Perturbation bounds of unitary and subunitary polar factors, \emph{SIAM J. Matrix Anal.\
Appl.} 23 (2002), 1183--1193 (as quoted in arXiv:2507.14940, Theorem~1.1).
\bibitem{candes} E.~J. Cand\`es, C.~A. Sing-Long and J.~D. Trzasko, Unbiased risk estimates for singular value thresholding
and spectral estimators, \emph{IEEE Trans.\ Signal Process.} 61 (2013).
\bibitem{cornish} E.~A. Cornish and R.~A. Fisher, Moments and cumulants in the specification of distributions, \emph{Rev.\
Inst.\ Internat.\ Statist.} 5 (1938).
\bibitem{minsker} S. Minsker and N. Strawn, The geometric median and applications to robust mean estimation,
arXiv:2307.03111.
\bibitem{schwank} R. Schwank and M. Drton, On the distance between mean and geometric median in high dimensions,
arXiv:2508.12926.
\bibitem{xie} X. Xie, Muon under gradient noise and the limits of orthogonalization near optima, arXiv:2609.32861.
\bibitem{chen} X. Chen, B. Li, L. Wang and P. Lu, When and how does Muon precondition? The role of structured gradient
noise, arXiv:2610.31728.
\bibitem{paquette} E. Paquette et al., Phases of Muon: when Muon eclipses SignSGD, arXiv:2605.09552.
\bibitem{ddj} T. Do, S. Dereich and A. Jentzen, On MUON optimization: from non-convergence to an error analysis with Polar
Express and the Newton--Schulz polynomial from implementations, arXiv:2608.04607.
\end{thebibliography}

\end{document}
